%\documentclass[12pt,oneside]{amsart}
\documentclass{scrartcl}

%\documentclass{article}

\usepackage{amsthm,amsmath,amssymb}
\usepackage[numbers]{natbib}

%\usepackage[colorlinks]{hyperref}
\usepackage{hyperref}
\hypersetup{
    colorlinks,%
    citecolor=black,%
    filecolor=black,%
    linkcolor=black,%
    urlcolor=black
}
\usepackage{hypernat}
\usepackage[top=2.5cm, bottom=2.5cm, left=2.8cm,
right=2.8cm]{geometry}
\usepackage{graphicx}


%\setlength{\parskip}{0pt}
%\setlength{\parsep}{0pt}
%\setlength{\headsep}{0pt}
%\setlength{\topskip}{0pt}
%\setlength{\topmargin}{0pt}
%\setlength{\topsep}{0pt}
%\setlength{\partopsep}{0pt}

\linespread{1}

%\usepackage{marvosym}

%\textwidth = 6.5 in \textheight = 9.2 in \oddsidemargin = 0.0 in
%\evensidemargin = 0.0 in \topmargin = -0.0 in \headheight = 0.0 in
%\headsep = 0.0in \parskip = 0.0in \parindent = 0.0in
%\def\baselinestretch{0.871}


\newtheorem{theorem}{Theorem}[section]
\newtheorem{proposition}[theorem]{Proposition}
\newtheorem{problem}[theorem]{Problem}
\newtheorem{fact}[theorem]{Fact}
\newtheorem{lemma}[theorem]{Lemma}
\newtheorem{defn}[theorem]{Definition}
\newtheorem{corollary}[theorem]{Corollary}
\newtheorem{remark}{Remark}[section]
\newtheorem{example}[theorem]{Example}

\newcommand{\E}{{\mathbb E}}
\newcommand{\se}{{\mathcal{E}}}
\newcommand{\R}{{\mathbb R}}
\renewcommand{\P}{{\mathbb P}}
\newcommand{\ac}{{\mathcal{A}}}
\newcommand{\A}{{\mathcal{A}}}
\newcommand{\B}{{\mathcal{B}}}
\newcommand{\C}{{\mathcal{C}}}
\newcommand{\M}{{\mathcal{M}}}
\newcommand{\Mf}{{\mathfrak{M}}}
\newcommand{\Qcal}{{\mathcal{Q}}}
\newcommand{\T}{{\mathcal{T}}}
\newcommand{\Z}{{\mathcal{Z}}}
\newcommand{\K}{{\mathcal{K}}}
\newcommand{\lc}{{\mathcal{L}}}
\newcommand{\Ps}{{\mathcal{P}}}
\newcommand{\G}{{\mathcal{G}}}
\newcommand{\pa}{{\mathring{p}}}
\newcommand{\F}{{\cal F}}
\newcommand{\samp}{{\mathcal{X}}}
\newcommand{\X}{{\mathcal{X}}}
\newcommand{\hi}{{\hat{\Phi}}}
\newcommand{\data}{{\text{data}}}
\newcommand{\relu}{{\text{ReLU}}}

\newcommand{\ridge}{{\hat{\beta}_{\text{ridge}}(\lambda)}}
\newcommand{\lasso}{{\hat{\beta}_{\text{lasso}}(\lambda)}}
\newcommand{\ridgemu}{{\hat{\mu}_t^{\text{ridge}}(\lambda)}}
\newcommand{\lassomu}{{\hat{\mu}_t^{\text{lasso}}(\lambda)}}


\newcounter{rcnt}[section]
\renewcommand{\thercnt}{(\roman{rcnt})}

\def\qt#1{\qquad\text{#1}}

\def\argmin{\mathop{\rm argmin}}
\def\argmax{\mathop{\rm argmax}}


\setlength{\parskip}{1.4 \medskipamount}
%\sloppy
%\linespread{1.3}

\begin{document}

\title{STAT 238 - Bayesian Statistics  \\
  Lecture Thirty Seven} 
\subtitle{Spring 2026, UC Berkeley}

\author{Aditya Guntuboyina}


\date{29 April 2026}

\maketitle

\section{Variational Inference}
We looked at the basics of variational inference in the last lecture. 
Consider
the standard Bayesian setup with data given by $y$ and parameters
given by $\theta$. The prior is $f_{\theta}(\theta)$ and the
likelihood is $f_{y \mid \theta}(y)$. The posterior is then:
\begin{align*}
  f_{\theta \mid y}(\theta) = \frac{f_{\theta}(\theta) f_{y \mid
  \theta}(y)}{f_y(y)}
\end{align*}
where the denominator is the evidence $f_{y}(y) = \int f_{y \mid
  \theta}(y) f_{\theta}(\theta) d\theta$.

The most important quantity in variational inference is the ELBO
defined as:
\begin{align}\label{elbo}
  \text{ELBO}(q) := \int q(\theta) \log \frac{f_{y, \theta}(y,
  \theta)}{q(\theta)} d\theta
\end{align}
where $q$ is an arbitrary probability density. The key facts about the
ELBO are:
\begin{align}\label{elbo_1}
\text{ELBO}(q) \leq \log f_{y}(y) ~~ \text{ for every } q.  
\end{align}
and
\begin{align}\label{elbo_2}
\max_{q} \text{ELBO}(q) = \log f_y(y) ~~ \text{  with the maximum
achieved at } q(\theta) = f_{\theta \mid y}(\theta).  
\end{align}
Both these facts are consequences of:
\begin{align*}
  \text{ELBO}(q) = \log f_{y}(y) - KL\left(q \| f_{\theta \mid y}
  \right). 
\end{align*}
Another useful fact about the ELBO (which follows directly from the
definition \eqref{elbo} when $f_{y, \theta}(y, \theta)$ is replaced by
$f_{y \mid \theta}(y) f_{\theta}(\theta)$) is:
\begin{align*}
  \text{ELBO}(q) = \int q(\theta) \log f_{y \mid \theta}(y) d\theta -
  KL(q \| f_{\theta}). 
\end{align*}

\section{Connection to the EM Algorithm}

The relation $\max_q \text{ELBO}(q) = \log f_{y}(y)$ (with the maximum
being achieved at the posterior $q(\theta) = f_{\theta \mid
  y}(\theta)$) can be used to understand the EM algorithm.

Consider the same Bayesian setting as in the previous section, but
suppose now that there is an additional hyperparameter $\alpha$, and
that both the likelihood $f_{y \mid \theta, \alpha}(y)$ as well as the
prior $f_{\theta \mid \alpha}(\theta)$ depend on $\alpha$.

The ELBO now becomes:
\begin{align*}
  \text{ELBO}(q, \alpha) = \int q(\theta) \log \frac{f_{y, \theta \mid
  \alpha}(y, \theta)}{q(\theta)} d\theta = \int q(\theta) \log
  \frac{f_{y \mid \theta, \alpha}(y) f_{\theta \mid
  \alpha}(\theta)}{q(\theta)} d\theta 
\end{align*}
and equation \eqref{elbo_2} now becomes:
\begin{align}\label{elbo_2_alpha}
  \max_{q} \text{ELBO}(q, \alpha) = \log f_{y\mid \alpha}(y) ~~ \text{  with the maximum
achieved at } q(\theta) = f_{\theta \mid y, \alpha}(\theta).  
\end{align}
Suppose now that we want to obtain the MLE $\hat{\alpha}$ of $\alpha$
i.e., we want to solve the optimization: 
\begin{align*}
\underset{\alpha}{\text{Maximize}}~ \log f_{y \mid \alpha}(y). 
\end{align*}
Using \eqref{elbo_2_alpha}, we can rewrite the optimization above as:
\begin{align*}
  \underset{\alpha}{\text{Maximize}}~ \max_{q} \text{ELBO}(q, \alpha), 
\end{align*}
which can be viewed as the two parameter maximization problem:
\begin{align*}
  \underset{q, \alpha}{\text{Maximize}} ~ \text{ELBO}(q, \alpha). 
\end{align*}
It is natural to solve this two-parameter optimization
alternatively. Fix $\alpha$ and optimize over $q$, and then fix $q$
and optimize over $\alpha$. This leads to the following algorithm.
\begin{enumerate}
\item Initialize $\alpha = \alpha^{(0)}$.
\item Repeat the following for $t = 0, 1, 2, \dots$ (until some  kind
  of convergence): 
  \begin{enumerate}
  \item Take $q^{(t)}$ to be the maximizer of $\text{ELBO}(q, \alpha^{(t)})$
    over all $q$ (for fixed $\alpha^{(t)}$)
  \item Take $\alpha^{(t+1)}$ to be the maximizer of $\text{ELBO}(q^{(t)},
    \alpha)$ over all $\alpha$. 
  \end{enumerate}
\end{enumerate}
By \eqref{elbo_2_alpha}, it is clear that $q^{(t)}$ which maximizes
$\text{ELBO}(q, \alpha^{(t)})$ is simply equal to:
\begin{equation*}
  q^{(t)} = f_{\theta \mid y, \alpha^{(t)}}. 
\end{equation*}
Thus the alternating minimization algorithm above is equivalent to:
\begin{enumerate}
\item Initialize $\alpha = \alpha^{(0)}$.
\item Repeat the following for $t = 0, 1, 2, \dots$ (until some  kind
  of convergence): 
  \begin{enumerate}
  \item Take $\alpha^{(t+1)}$ to be the maximizer of
    $\text{ELBO}(f_{\theta, \mid y, \alpha^{(t)}}, 
    \alpha)$ over all $\alpha$. 
  \end{enumerate}
\end{enumerate}
The quantity $\text{ELBO}(f_{\theta, \mid y, \alpha^{(t)}},
\alpha)$ equals:
\begin{align*}
  \text{ELBO}(f_{\theta, \mid y, \alpha^{(t)}},
\alpha) &= \int f_{\theta \mid y, \alpha^{(t)}}(\theta) \log
          \frac{f_{y, \theta \mid \alpha}(y, \theta)}{f_{\theta \mid
          y, \alpha^{(t)}}(\theta)} d\theta  \\
  &= \int f_{\theta \mid y, \alpha^{(t)}}(\theta) \log f_{y, \theta
    \mid \alpha}(\theta) d\theta - \int f_{\theta \mid y,
    \alpha^{(t)}}(\theta) \log f_{\theta \mid y, \alpha^{(t)}}(\theta)
    d\theta. 
\end{align*}
The second term actually does not depend on $\alpha$, so maximizing
$\text{ELBO}(f_{\theta, \mid y, \alpha^{(t)}}, \alpha)$  over $\alpha$
is equivalent to maximizing the first term above:
\begin{align}\label{Et}
  E(t, \alpha) := \int f_{\theta \mid y, \alpha^{(t)}}(\theta) \log f_{y, \theta
    \mid \alpha}(\theta) d\theta. 
\end{align}
The alternating maximization algorithm above then becomes:
\begin{enumerate}
\item Initialize $\alpha = \alpha^{(0)}$.
\item Repeat the following for $t = 0, 1, 2, \dots$ (until some  kind
  of convergence): 
  \begin{enumerate}
  \item Calculate $E(t, \alpha)$ given in \eqref{Et}
  \item Take $\alpha^{(t+1)}$ to be the maximizer of
    $E(t, \alpha)$ over all $\alpha$. 
  \end{enumerate}
\end{enumerate}
This is the EM algorithm. The calculation of $E(t, \alpha)$ is called
the $E$-step and the maximization of $E(t, \alpha)$ over $\alpha$ is
called the $M$-step. For more on this connection between Variational
Inference (VI) and EM,
see \citet{neal1998view}. 

\section{Coordinatewise Variational Inference}

The key problem that we need to solve in VI is that of maximizing
$\text{ELBO}(q)$. The full maximizer of $\text{ELBO}(q)$ over $q$is, 
as already mentioned, $f_{\theta \mid y}$. But this posterior
is intractable in general and the very reason for using variational
inerference is this intractability. The strategy is to restrict the class
of $q$ over which we look for minimizers. This hopefully will lead to
a \textit{tractable} minimizer which is close to the actual posterior
$f_{\theta \mid y}$.

Suppose $\theta$ can be decomposed into $k$ separate classes of
parameters as $\theta = (\theta_1, \dots, \theta_k)$. We look for
variational approximations of the form $q(\theta) = q_1(\theta_1)
\dots q_k(\theta_k)$. The ELBO maximization problem is:
\begin{equation}\label{kfe}
  \underset{q_1, \dots, q_k}{\text{maximize}}~ \text{ELBO}(q_1, \dots,
  q_k) = \int \dots \int 
  q_1(\theta_1) \dots q_k(\theta_k) \log \frac{f_{y, \theta}(y, \theta)}{q_1(\theta_1) \dots
    q_k(\theta_k)} d\theta.    
\end{equation}
Alternating minimization leads to the following result. 

\begin{proposition}[CAVI]\label{cavi}
  Consider the maximization problem \eqref{kfe}. Fix $j = 1,
  \dots, k$ and fix $q_i, i \neq j$. Then the minimizing density $q_j
  = q_j^*$ is given by
  \begin{equation}
    \label{cavi_j}
    q_j^*(\theta_j) \propto \exp \left[\int q_{-j}(\theta_{-j}) \log
      f_{\theta_j \mid y, \theta_{-j}, \alpha}(\theta_j) d\theta_{-j}
    \right]. 
  \end{equation}
  where $q_{-j}(\theta_j)$ denotes the product of $q_{l}(\theta_l)$
  over all $l \neq j$ and $d\theta_j$ is the product of $d\theta_l$
  over all $l \neq j$. 
\end{proposition}
The above result is a consequence of the following simple fact.

\begin{lemma}\label{klfacts}
  \begin{enumerate}
  \item Suppose $g$ is a nonnegative function that is not necessarily
    a density function. Then the minimizer of 
    \begin{equation*}
 L(f) := \int f \log \frac{f}{g}
\end{equation*}
over all densities $f$ is given  by
\begin{equation}\label{klfact1}
  f^*(x) \propto g(x) 
\end{equation}
which of course means that $f^* = g/(\int g)$. 
  \item Suppose $f_2(\cdot)$ below is a fixed probability density
    function, and $g(x_1, x_2)$ is a fixed nonnegative function that
    is not 
    necessarily a density. Then the minimizer of 
    \begin{equation*}
\Gamma(f_1) := \int \int f_1(x_1) f_2(x_2) \log
\frac{f_1(x_1)f_2(x_2)}{g(x_1, x_2)} dx_1 dx_2 
    \end{equation*}
    over all densities $f_1$ is given by
    \begin{equation}\label{klfact2}
      f_1^*(x_1) \propto \exp \left( \int f_2(x_2) \log g(x_1, x_2)
        dx_2 
      \right). 
    \end{equation}
  \end{enumerate}
\end{lemma}

\begin{proof}[Proof of Lemma \ref{klfacts}]
  The first fact is easy because we can write
  \begin{align*}
    L(f) &= \int f \log \frac{f}{g} \\ &= \int f \log \frac{f}{(\int
      g)^{-1} g} - \log (\int g) \int f \\ &= \int f \log \frac{f}{(\int
      g)^{-1} g} - \log \int g = KL(f \| (\int g)^{-1} g) - \log \int
                                         g, 
  \end{align*}
  so it is clear that the minimizing $f$ equals $(\int g)^{-1} g$.

  For the second fact, write
  \begin{align*}
    \Gamma(f_1) &= \int \int f_1(x_1) f_2(x_2) \log
                  \frac{f_1(x_1)f_2(x_2)}{g(x_1, x_2)} dx_1 dx_2  \\
    &= \int \int f_1(x_1) f_2(x_2) \left[\log
      \frac{f_1(x_1)}{g(x_1, x_2)} + \log f_2(x_2) \right] dx_1 dx_2 \\
    &= \int \int f_1(x_1) f_2(x_2) \log
      \frac{f_1(x_1)}{g(x_1, x_2)} + \int f_2(x_2) \log f_2(x_2)
      dx_2. 
  \end{align*}
  We can ignore the second term as it does not depend on $f_1$. So
  minimizing $\Gamma(f_1)$ is the same as minimizing the first term: 
  \begin{align*}
    &  \int \int f_1(x_1) f_2(x_2) \log \frac{f_1(x_1)}{g(x_1, x_2)}
      dx_1 d_2 \\
    &= \int f_1(x_1) \int f_2(x_2) \left[\log f_1(x_1) - \log g(x_1,
      x_2) \right] dx_2 dx_1 \\
    &= \int f_1(x_1) \left[\log f_1(x_1) - \int f(x_2) \log g(x_1,
      x_2) dx_2 \right] dx_1 \\
    &= \int f_1(x_1) \log \frac{f_1(x_1)}{\exp \left[\int f(x_2) \log
      g(x_1, x_2) dx_2 \right]} dx_1. 
  \end{align*}
  We can thus use the first fact with $g(x_1) = \exp \left[ \int f(x_2) \log
      g(x_1, x_2) dx_2\right]$ to claim that the minimizer $f_1^*$ of
    $\Gamma(f_1)$ satisfies \eqref{klfact2}. 
\end{proof}







\begin{thebibliography}{99}

\bibitem[Neal and Hinton(1998)]{neal1998view}
Neal, R.~M. and Hinton, G.~E. (1998).
\newblock A view of the EM algorithm that justifies incremental, sparse, and other variants.
\newblock In M.~I. Jordan (ed.), \emph{Learning in Graphical Models},
pages 355--368. Springer Netherlands.  

%  \bibitem[Neal(2011)]{neal2011mcmc}
%Radford M. Neal.
%\newblock MCMC using Hamiltonian dynamics.
%\newblock In \emph{Handbook of Markov Chain Monte Carlo}, pages 47--95.
%\newblock Chapman and Hall/CRC, 2011.

%\bibitem[Girolami and Calderhead(2011)]{girolami2011riemann}
%Mark Girolami and Ben Calderhead.
%\newblock Riemann manifold Langevin and Hamiltonian Monte Carlo methods.
%\newblock \emph{Journal of the Royal Statistical Society: Series B (Statistical Methodology)}, 73(2):123--214, 2011.

%  \bibitem[Lai et~al.(2025)Lai, Song, Kim, Mitsufuji, and Ermon]{lai2025principles}
%Lai, C.-H., Song, Y., Kim, D., Mitsufuji, Y., and Ermon, S. (2025).
%\newblock The principles of diffusion models.
%\newblock \emph{arXiv preprint arXiv:2510.21890}.

%\bibitem[Glatt-Holtz et~al.(2024)]{glatt2024sacred}
%Nathan E. Glatt-Holtz, Andrew J. Holbrook, Justin A. Krometis, Cecilia F. Mondaini, and Ami Sheth.
%\newblock Sacred and Profane: from the Involutive Theory of MCMC to Helpful Hamiltonian Hacks.
%\newblock \emph{arXiv preprint arXiv:2410.17398}, 2024.

%\bibitem[Roberts and Rosenthal(2004)]{RobertsRosenthal2004}
%Roberts, G. O. and Rosenthal, J. S. (2004).
%General state space Markov chains and MCMC algorithms.
%\textit{Probability Surveys}, \textbf{1}, 20--71.

%\bibitem[Meyn and Tweedie(2009)]{MeynTweedie2009}
%Meyn, S. P. and Tweedie, R. L. (2009).
%\textit{Markov Chains and Stochastic Stability}, 2nd ed.
%Cambridge University Press.


  
%\bibitem[Chib and Greenberg(1995)]{chib1995}
%Chib, S. and Greenberg, E. (1995).
%Understanding the Metropolis--Hastings algorithm.
%\textit{The American Statistician}, \textbf{49}, 327--335.


%\bibitem[MacKay and Peto(1995)]{MacKay1995HierarchicalDirichlet}
%D.~J.~C. MacKay and L.~C. Peto,
%\newblock ``A hierarchical Dirichlet language model,''
%\newblock \emph{Natural Language Engineering}, vol.~1,
%no.~3, pp.~289--308, 1995.

%\bibitem[Rubin(1981)]{Rubin1981BayesianBootstrap}
%D.~B. Rubin,
%\newblock ``The Bayesian bootstrap,''
%\newblock \emph{The Annals of Statistics}, vol.~9,
%no.~1, pp.~130--134, 1981.
  
%\bibitem[Fisher(1934)]{Fisher1934}
%R.~A. Fisher,
%\newblock ``Probability, likelihood and quantity of information in the logic of
%uncertain inference,''
%\newblock \emph{Proceedings of the Royal Society of London. Series~A,
%Containing Papers of a Mathematical and Physical Character}, vol.~146,
%no.~856, pp.~1--8, 1934.

%\bibitem[Efron(2010)]{Efron}
%Efron, B. (2010)
%\textit{Large-scale inference: empirical Bayes methods for estimation,
%testing, and prediction}. Cambridge University Press, 2010.

%\bibitem[Shumway and Stoffer(2017)]{ShumwayStoffer}
%Shumway, R. H. and Stoffer, D. S. (2017)
%\textit{Time Series Analysis and Its Applications: With R
%  Examples}. Springer, 4th edition. 
  
%     \bibitem[Jaynes(2003)]{Jaynes} Jaynes, E. T, (2003)
%  \textit{Probability theory: the logic of science}. Cambridge
%  University Press, 2003.
 
%  \bibitem[Sinai(1992)]{Sinai} Sinai, Y. G, (1992)
%  \textit{Probability theory: an introductory course}. Springer
%  Verlag, 1992.  

%\bibitem[{deGroot(1986)}]{DeGrootBlackwell}DeGroot, M. H. (1986)
%       A conversation with David Blackwell.
%\newblock \emph{Statistical Science}, \textbf{1}, 40--53.

%         \bibitem[Mosteller(1987)]{Mosteller} Mosteller, F, (1987)
%  \textit{Fifty challenging problems in probability with
%    solutions}. Courier Corporation, 1987.

%    \bibitem[MacKay(2003)]{MacKay} MacKay, D. J. C., (2003)
%  \textit{Information theory, inference, and learning
%  algorithms}. Cambridge University Press, 2003.

%\bibitem[{Lindley(2013)}]{Lindley}Lindley, D. V. (2013)
%   \textit{Understanding Uncertainty}. John Wiley and sons, 2013.


\end{thebibliography}







\end{document}

%%% Local Variables:
%%% mode: latex
%%% TeX-master: t
%%% End:
